% fancybeamer.tex — mdOffice Fancy Beamer style
% Loaded via beamer-style: "fancybeamer" in the document front-matter.
%
% ---- FRONTMATTER EXAMPLE ---------------------------------------------------
%
% ---
% make-beamer:  true
% beamer-style: "fancybeamer"
%
% title:        "Lecture Title"
% author:       "Your Name"
% date:         "Date"
%
% aspectratio:  169          # 16:9 (default)
% fontsize:     11pt
% sansfont:     "Fira Sans"  # font family — any system font name (XeLaTeX)
% ---
%
% ---- SLIDE HEADER COLORS ---------------------------------------------------
% Place one of these as the FIRST line of a slide's content:
%   \myheader   \blueheader   \redheader   \greenheader   \cyanheader   \grayheader
% The default is MyHeader — slides without a macro inherit the previous color.
%
% ---- THEOREM BOXES ---------------------------------------------------------
% Numbered, with title and cross-reference label:
%   \begin{bluebox}{Title}{label} ... \end{bluebox}
%   \begin{redbox}{Title}{label}  ... \end{redbox}
%   \begin{greenbox}{Title}{label}... \end{greenbox}
%   \begin{cyanbox}{Title}{label} ... \end{cyanbox}
%   \begin{graybox}{Title}{label} ... \end{graybox}
% Starred form (no counter):
%   \begin{bluebox*}{Title} ... \end{bluebox*}
%
% ---- FOCUS BOXES -----------------------------------------------------------
% No title, lightly tinted background — for key formulas, remarks, warnings:
%   \begin{myfocus}   ... \end{myfocus}
%   \begin{bluefocus} ... \end{bluefocus}
%   \begin{redfocus}  ... \end{redfocus}
%   \begin{greenfocus}... \end{greenfocus}
%   \begin{cyanfocus} ... \end{cyanfocus}
%   \begin{grayfocus} ... \end{grayfocus}
%
% ---- WARNING / INFO MACROS -------------------------------------------------
%   \warn                              — red triangle + "Not part of the syllabus"
%   \warnCustomMsg[color]{text}        — custom triangle message
%   \info                              — blue circle  + "For your information"
%   \infoCustomMsg[color]{text}        — custom info  message
%
% ---- MATH MACROS -----------------------------------------------------------
%   \norm{v}   →  ‖v‖        \abs{x}   →  |x|
% ----------------------------------------------------------------------------


% ---------------------------------------------------------------
% BASE THEME — default gives a clean slate for full customization
\usetheme{default}
% ---------------------------------------------------------------


% ---------------------------------------------------------------
% PACKAGES
\usepackage[most]{tcolorbox}   % theorem/focus boxes
\usepackage{fontawesome5}       % icons for \warn / \info
\usepackage{xspace}             % smart spacing after macros
% ---------------------------------------------------------------


% ---------------------------------------------------------------
% FONT — set via front-matter: sansfont: "Font Name"
% pandoc injects \setsansfont{...} automatically for XeLaTeX.
% No font set here — falls back to XeLaTeX's default sans-serif.
%
% Popular choices (must be installed as system fonts):
%   Fira Sans, Open Sans, Source Sans Pro      (Google Fonts)
%   Arial, Calibri, Trebuchet MS               (Windows)
%   DejaVu Sans, Liberation Sans               (Linux / cross-platform)
% ---------------------------------------------------------------


% ---------------------------------------------------------------
% COLORS
%
% 1. Base colors — raw RGB palette, used for inline text and as source for derived colors
\definecolor{myblue}{RGB}{30, 90, 160}
\definecolor{myred}{RGB}{180, 40, 40}
\definecolor{mygreen}{RGB}{46, 125, 50}
\definecolor{mycyan}{RGB}{0, 170, 180}
\definecolor{mygray}{RGB}{64, 64, 64}
%
% 2. Heading colors — for slide headers, block titles, theorem box frames
%    Adjust the ! expression to darken/shift without touching the base RGB
\colorlet{myblueheading}{myblue!90}
\colorlet{myredheading}{myred!90}
\colorlet{mygreenheading}{mygreen!90}
\colorlet{mycyanheading}{mycyan!90}
\colorlet{mygrayheading}{mygray!90}
%
% 3. Focus colors — light tints of the base colors
%    Used for focus box backgrounds and block body backgrounds
\colorlet{mybluefocus}{myblue!5}
\colorlet{myredfocus}{myred!5}
\colorlet{mygreenfocus}{mygreen!5}
\colorlet{mycyanfocus}{mycyan!5}
\colorlet{mygrayfocus}{mygray!5}
%
% Header expressions — \colorbox evaluates ! tints at render time
\def\HeaderDefault{myblueheading}
\def\HeaderBlue{myblueheading}
\def\HeaderRed{myredheading}
\def\HeaderGreen{mygreenheading}
\def\HeaderCyan{mycyanheading}
\def\HeaderGray{mygrayheading}
% ---------------------------------------------------------------

% ---------------------------------------------------------------
% STRUCTURAL COLOR (navigation dots, sidebar elements, etc.)
\setbeamercolor{structure}{fg=myblue}
% ---------------------------------------------------------------

% ---------------------------------------------------------------
% STANDARD BLOCK COLORS (used by pandoc-generated block/alertblock/exampleblock)
\setbeamercolor{block title}{fg=white, bg=myblueheading}
\setbeamercolor{block body}{bg=mybluefocus}
\setbeamercolor{block title alerted}{fg=white, bg=myredheading}
\setbeamercolor{block body alerted}{bg=myredfocus}
\setbeamercolor{block title example}{fg=white, bg=mygreenheading}
\setbeamercolor{block body example}{bg=mygreenfocus}
% ---------------------------------------------------------------


% ---------------------------------------------------------------
% BEAMER TEMPLATE OVERRIDES
\setbeamercovered{transparent}
\setbeamertemplate{headline}{}
\setbeamertemplate{navigation symbols}{}

% Optional bottom-right logo (from front-matter: logo + logo-height)
% Apply at begin-document so this wins over pandoc's default \logo metadata.
\ifdefined\logopath
  \ifdefined\logoheight\else\def\logoheight{0.8cm}\fi
  \AtBeginDocument{\logo{\includegraphics[height=\logoheight]{\logopath}}}
\fi

% Custom bullet points
\setbeamertemplate{itemize item}{\textcolor{black}{$\bullet$}}
\setbeamertemplate{itemize subitem}{\textcolor{black}{$\bullet\,\bullet$}}
\setbeamertemplate{itemize subsubitem}{\textcolor{black}{$\bullet\bullet\bullet$}}

% Footline: frame number only
\setbeamertemplate{footline}[frame number]

% ---------------------------------------------------------------
% FRAME HEADER COLOR SYSTEM
%
% \fancyheadcolor stores the active color name (e.g. "HeaderBlue").
% \gdef is used everywhere so the value is truly global and survives
% beamer's per-frame TeX groups.
%
% The frametitle template reads \fancyheadcolor at render time (which
% happens after \end{frame} closes the frame group), so a macro placed
% anywhere in a slide's body colors THAT slide's title bar.
%
% Recommended pattern — put the macro as the FIRST line of a slide:
%
%   ## My Slide
%   \redheader
%   Content…
%
% The color persists to subsequent slides until another macro overrides it.
\gdef\fancyheadcolor{\HeaderDefault}

\newcommand{\blueheader}{\gdef\fancyheadcolor{\HeaderBlue}}
\newcommand{\redheader}{\gdef\fancyheadcolor{\HeaderRed}}
\newcommand{\greenheader}{\gdef\fancyheadcolor{\HeaderGreen}}
\newcommand{\cyanheader}{\gdef\fancyheadcolor{\HeaderCyan}}
\newcommand{\grayheader}{\gdef\fancyheadcolor{\HeaderGray}}
\newcommand{\myheader}{\gdef\fancyheadcolor{\HeaderDefault}}

% Frametitle: full-width colored bar.
% \fancyheadcolor always expands to a pre-defined named color (myredheading etc.),
% so \setbeamercolor can accept it and beamercolorbox fills the full paper width.
\setbeamertemplate{frametitle}{%
  \nointerlineskip
  \edef\fancybeamer@bgcol{\fancyheadcolor}%
  \setbeamercolor{fancybeamer@hdr}{bg=\fancybeamer@bgcol, fg=white}%
  \begin{beamercolorbox}[wd=\paperwidth, ht=2.1ex, dp=0.8ex, leftskip=0.3cm]{fancybeamer@hdr}%
    \usebeamerfont{frametitle}\insertframetitle%
  \end{beamercolorbox}%
}
% ---------------------------------------------------------------


% ---------------------------------------------------------------
% TCOLORBOX GLOBAL DEFAULTS
\tcbset{
  colback=white,
  fonttitle=\bfseries,
  breakable,
  enhanced,
  top=2pt, bottom=2pt, left=2pt, right=2pt,
  boxsep=3pt
}

% Theorem-style boxes: \begin{bluebox}{Title}{label} ... \end{bluebox}
% All share a single counter so numbering is consistent across colors.
\newtcbtheorem{bluebox}{}%
  {colframe=myblueheading}{bluebox}

\newtcbtheorem[use counter from=bluebox]{redbox}{}%
  {colframe=myredheading}{redbox}

\newtcbtheorem[use counter from=bluebox]{greenbox}{}%
  {colframe=mygreenheading}{greenbox}

\newtcbtheorem[use counter from=bluebox]{cyanbox}{}%
  {colframe=mycyanheading}{cyanbox}

\newtcbtheorem[use counter from=bluebox]{graybox}{}%
  {colframe=mygrayheading}{graybox}

% Focus boxes: no title, lightly tinted background
\newtcolorbox{myfocus}{colback=mybluefocus,   colframe=white}
\newtcolorbox{bluefocus}{colback=mybluefocus, colframe=white}
\newtcolorbox{redfocus}{colback=myredfocus,   colframe=white}
\newtcolorbox{greenfocus}{colback=mygreenfocus, colframe=white}
\newtcolorbox{cyanfocus}{colback=mycyanfocus, colframe=white}
\newtcolorbox{grayfocus}{colback=mygrayfocus, colframe=white}
% ---------------------------------------------------------------


% ---------------------------------------------------------------
% WARNING AND INFO MACROS
\newcommand{\warn}{%
  \textcolor{red}{\faExclamationTriangle\ \textbf{Not part of the syllabus}}\xspace%
}
\newcommand{\warnCustomMsg}[2][red]{%
  \textcolor{#1}{\faExclamationTriangle\ \textbf{#2}}\xspace%
}
\newcommand{\info}{%
  \textcolor{blue}{\faInfoCircle\ \textbf{For your information}}\xspace%
}
\newcommand{\infoCustomMsg}[2][blue]{%
  \textcolor{#1}{\faInfoCircle\ \textbf{#2}}\xspace%
}
% ---------------------------------------------------------------


% ---------------------------------------------------------------
% MATH MACROS
% \providecommand avoids "already defined" errors if the user loads
% a package such as physics or mathtools that defines these names.
\providecommand{\norm}[1]{\left\lVert#1\right\rVert}
\providecommand{\abs}[1]{\left\lvert#1\right\rvert}
% ---------------------------------------------------------------
